\begin{lemma}[The actual zero-scheme dimension is the original line degree]
\label{lem:fc-original-zero-section-degree-dimension}
\leanok
\lean{EnvelopeProperty.CurveLineZeroSections.degree_eq_zeroScheme_finrank, EnvelopeProperty.CurveLineZeroSections.degree_pos_of_nonempty_zeroScheme}
\uses{thm:fc-original-nonzero-section-degree-nonnegative}
Let $X$ be the original integral projective curve over $k$, of dimension
at most one.  For any original invertible $L$ and nonzero original
$s:\mathcal O_X\to L$, let $Z(s)$ be its actual zero subscheme.  Then
\[
 \deg_k L=\dim_k\Gamma(Z(s),\mathcal O_{Z(s)}).
\]
If $Z(s)$ is nonempty, then $\deg_k L>0$.
\end{lemma}
\begin{proof}
The actual dual evaluation of $s$ identifies the whole $L^\vee$ with its
actual Cartier image ideal $I$.  Its support is proper and its actual
closed subscheme is finite over $k$.  The full coherent Euler exact
sequence gives $\deg I=-\dim_k\Gamma(Z(s),\mathcal O_{Z(s)})$.
The true tensor/dual degree formula gives $\deg L^\vee=-\deg L$,
yielding the asserted equality.  If $Z(s)$ has a point then its entire
section ring is nontrivial.  This ring, with the actual original
coefficient action, is finite dimensional over $k$, so its dimension is
strictly positive.
\end{proof}

\begin{lemma}[Empty actual zero ideal gives a whole line trivialization]
\label{lem:fc-original-zero-section-whole-trivialization}
\leanok
\lean{EnvelopeProperty.CurveLineZeroSections.topIdealModuleIso, EnvelopeProperty.CurveLineZeroSections.trivializationOfZeroIdealTop}
\uses{lem:fc-original-zero-section-degree-dimension}
For any original integral scheme, the actual unit ideal module is the
structure module under its original inclusion.  If an original nonzero
section of an invertible $L$ has unit zero ideal, it internally gives a
whole isomorphism $L\simeq\mathcal O_X$.
\end{lemma}
\begin{proof}
Every section of the structure module belongs to the unit ideal on every
affine subopen, so the original whole ideal inclusion is an isomorphism.
Compose this with the actual dual-section/image isomorphism to obtain
$L^\vee\simeq\mathcal O_X$.  Tensor its inverse with $L$ and use the
original canonical evaluation $L\otimes L^\vee\simeq\mathcal O_X$.
This yields the whole original line isomorphism.
\end{proof}

\begin{theorem}[Globally generated degree zero detects the whole trivial line]
\label{thm:fc-original-generated-line-degree-zero-trivial}
\leanok
\lean{EnvelopeProperty.CurveLineZeroSections.zeroIdeal_eq_top_of_degree_eq_zero, EnvelopeProperty.CurveLineZeroSections.degree_eq_zero_iff_trivial_of_generatingSections, EnvelopeProperty.CurveLineZeroSections.degree_pos_of_generatingSections_of_nontrivial}
\uses{lem:fc-original-zero-section-degree-dimension,lem:fc-original-zero-section-whole-trivialization}
On the same original integral projective curve, a genuinely globally
generated original invertible line has degree zero if and only if the
whole original line is trivial.  Consequently every such nontrivial
original line has strictly positive degree.
\end{theorem}
\begin{proof}
Actual generation at the original generic point supplies a nonzero
original generator.  If its actual zero subscheme were nonempty, the
preceding formula would give strictly positive degree.  Thus degree zero
forces empty actual zero support, hence the unit zero ideal, and the
preceding whole trivialization applies.  Conversely an actual whole
isomorphism to the structure line gives degree zero by invariance of the
true full Euler degree.  The nonnegative degree already proved for
globally generated lines is therefore strictly positive in the
nontrivial case.
\end{proof}
